\appendix
\chapter{{Experimentelle Studie 1 -- Ansicht eines Items}} \label{Screen}
\begin{figure}
  \includegraphics[angle=-90, width=0.67\textwidth]{Screen}
  \caption[Exemplarische Ansicht eines Items von Experiment 1] {Exemplarische Ansicht eines Items von Experiment 1.}
  \label{Screen1}
\end{figure}

\chapter{{Experimentelle Studie 1 -- Ratings der Intensivierer}} \label{Rat1}
\vspace*{-\baselineskip}
\begin{figure}[h!]
\includegraphics[angle=-90, width=0.8\textwidth]{Skalar7_neu}
\caption[Säulendiagramm zu den einzelnen Intensivierern bei Aufgabe I von Experiment 1]{Säulendiagramm zu den einzelnen Intensivierern bei Aufgabe I -- Skalarität ({x-Achse =} Intensivierer, y-Achse = Skalenwert).}
\label{Skalar7}
\end{figure}

\chapter{{Experimentelle Studie 1 -- Ratings der Pseudowörter}}  \label{Pseu1}
\begin{table}
\begin{tabularx}{.8\textwidth}{lYY}
\lsptoprule
   \multicolumn{3}{c}{\textbf{Aufgabe I Skalarität}}   \\
   \midrule
{Pseudonomen -- Pseudoadjektiv} & Skalenwert  &   \textit{SD} \\
\midrule
 \rowcolor[gray]{.90}  \textit{Der Gedalst -- wübig} &   62,55   & 25,0  \\
                        \textit{Das Zülugs -- wosplich} &  62,76    & 24,0 \\
  \rowcolor[gray]{.90}  \textit{Das Genavs -- laukrig} &  63,39    & 24,3 \\
                        \textit{Das Zeseet -- bühllich} &   63,76    & 24,4  \\
    \rowcolor[gray]{.90} \textit{Der Cobe -- speiglich} &  63,79    & 24,1 \\
                        \textit{Der Wonger -- zänsig} &  64,25    & 24,8  \\
    \rowcolor[gray]{.90} \textit{Der Zegeff -- ralklich} &  64,81    & 25,4   \\
                        \textit{Der Zotzer -- galklich} &  65,45    & 21,7  \\
   \rowcolor[gray]{.90}   \textit{Die Itne -- schühnlich} &  65,48  & 23,1   \\
                          \textit{Das Tusern -- zarwig} &  65,54    & 24,5  \\
   \rowcolor[gray]{.90}   \textit{Der Pulpart -- röngig} & 65,55   & 24,4  \\
                           \textit{Der Nümig -- fongig} &  65,94    & 25,0  \\
    \rowcolor[gray]{.90}     \textit{Der Conen -- waumlich} &  66,08 & 23,7   \\
                             \textit{Der Jordruck -- kieflich} &  67,02    & 23,8  \\
      \rowcolor[gray]{.90}  \textit{Das Hensuln -- ontig} &  67,57  & 22,1  \\
                            \textit{Die Itse -- rümmig} &  69,05   & 21,6  \\
     \rowcolor[gray]{.90}   \textit{Der Darster -- ormig} &  69,63    & 22,8  \\
                          \textit{Das Zafche -- göcktig} &  69,69    & 22,0  \\
\lspbottomrule
\end{tabularx}
\caption[Gemittelte Skalenwerte der Pseudowörter bei {Aufgabe I} von Experiment 1]{Gemittelte Skalenwerte und Standardabweichungen der Kombination aus Pseudonomen und Pseudoadjektiv bei Aufgabe I -- Skalarität (angeordnet nach ansteigendem Skalenwert).}
\label{Pseudos1}
\end{table}

\chapter{{Experimentelle Studie 1 -- Ratings der Pseudowörter}}  \label{Happy0}
\vspace{-\baselineskip}
\begin{figure}[h!]
\includegraphics[angle=-90, width=0.8\textwidth]{Happy1}
\caption[Säulendiagramm zu den einzelnen Pseudowörtern bei Aufgabe I von Experiment 1]{Säulendiagramm zu den einzelnen Pseudowort-Kombinationen bei Aufgabe I -- Skalarität ({x-Achse =} Pseudowörter, y-Achse = Skalenwert).}
\label{Happy3}
\end{figure}

\chapter{{Experimentelle Studie 1 -- Ratings der Intensivierer}}  \label{Rat2}
\vspace{-\baselineskip}
\begin{figure}[h!]
\includegraphics[angle=-90, width=0.8\textwidth]{Expressiv7_neu}
\caption[Säulendiagramm zu den einzelnen Intensivierern bei Aufgabe II von Experiment 1]{Säulendiagramm zu den einzelnen Intensivierern bei Aufgabe II -- Sprechereinstellung ({x-Achse =} Intensivierer, y-Achse = Skalenwert).}
\label{Expressiv7}
\end{figure}

\chapter{{Experimentelle Studie 1 -- Ratings der Pseudowörter}}   \label{Pseu2}

\begin{table}
\begin{tabularx}{.8\textwidth}{lYY}
\lsptoprule
   \multicolumn{3}{c}{\textbf{Aufgabe II Sprechereinstellung}}   \\
   \midrule
{Pseudonomen -- Pseudoadjektiv} & Skalenwert  &   \textit{SD} \\
\midrule
 \rowcolor[gray]{.90}     \textit{Der Cobe -- speiglich} &  58,75   & 27,6  \\
                        \textit{Der Wonger -- zänsig} &  59,12    & 25,4  \\
  \rowcolor[gray]{.90}    \textit{Das Tusern -- zarwig} &  60,11   & 27,5   \\
                        \textit{Der Gedalst -- wübig} &   60,25   & 27,1  \\
    \rowcolor[gray]{.90}    \textit{Der Jordruck -- kieflich} &  60,81  & 27,8  \\
                        \textit{Der Nümig -- fongig} &  60,84    & 26,7  \\
 \rowcolor[gray]{.90}     \textit{Das Zülugs -- wosplich} &  61,16  & 27,6    \\
                      \textit{Der Pulpart -- röngig} & 61,37    & 24,7  \\
    \rowcolor[gray]{.90}     \textit{Der Conen -- waumlich} &  61,39    & 26,8  \\
                        \textit{Das Zeseet -- bühllich} &   61,96    & 25,9  \\
    \rowcolor[gray]{.90}   \textit{Der Darster -- ormig} &  62,66   & 28,6  \\
                        \textit{Der Zotzer -- galklich} &  62,66    & 25,4  \\
   \rowcolor[gray]{.90}    \textit{Die Itse -- rümmig} &  62,76  & 26,2  \\
                      \textit{Das Zafche -- göcktig} &  63,88    & 25,3  \\
    \rowcolor[gray]{.90}   \textit{Das Hensuln -- ontig} &  64,15    & 29,6 \\
                        \textit{Das Genavs -- laukrig} &  64,48   & 27,9   \\
   \rowcolor[gray]{.90}   \textit{Der Zegeff -- ralklich} &  64,69   & 25,8    \\
                        \textit{Die Itne -- schühnlich} &  65,34     & 24,3  \\
\lspbottomrule
\end{tabularx}
\caption[Gemittelte Skalenwerte der Pseudowörter bei {Aufgabe II} von Experiment 1]{Gemittelte Skalenwerte und Standardabweichungen der Kombination aus Pseudonomen und Pseudoadjektiv bei Aufgabe II -- Sprechereinstellung (angeordnet nach ansteigendem Skalenwert).}
\label{Pseudos2}
\end{table}

\chapter{{Experimentelle Studie 1 - Ratings der Pseudowörter}}  \label{Happy4}
\vspace*{-\baselineskip}

\begin{figure}[h!]
\includegraphics[angle=-90, width=0.8\textwidth]{Happy3}
\caption[Säulendiagramm zu den einzelnen Pseudowörtern bei Aufgabe II von Experiment 1]{Säulendiagramm zu den einzelnen Pseudowort-Konbinationen bei Aufgabe II -- Sprechereinstellung ({x-Achse =} Pseudowörter, {y-Achse =} Skalenwert).}
\label{Happy5}
\end{figure}

\chapter{{Experimentelle Studie 1 -- Ratings der Pseudowörter}}  \label{Pseu3}

\begin{table}
\begin{tabularx}{\textwidth} {lYYYYYY}

\lsptoprule
 & \multicolumn{6}{c}{\textbf{Aufgabe III Pseudowörter}}\\
 \midrule
 & \multicolumn{6}{c}{\textbf{Einstellungspolarität}}  \\
  &\multicolumn{2}{c}{\textsc{positiv}} & \multicolumn{2}{c}{\textsc{neutral}} & \multicolumn{2}{c}{\textsc{negativ}} \\
 \cmidrule(lr){2-3}\cmidrule(lr){4-5}\cmidrule(lr){6-7}
  & Taler& \textit{SD}  & Taler& \textit{SD}  & Taler& \textit{SD}   \\
\midrule
 \rowcolor[gray]{.90}   \textit{Der Conen -- waumlich} & 45,45& 32,3 &  24,37 & 23,9 & 30,18  & 29,7\\
                  \textit{Der Darster -- ormig} & 30,06 & 28,4 & 34,78& 30,9  & 35,16 & 33,6    \\
 \rowcolor[gray]{.90}    \textit{Die Itne -- schühnlich} & 33,87& 29,0 & 30,73  & 29,4 & 35,40 & 31,0  \\
                        \textit{Das Genavs -- laukrig} & 31,12& 28,3 & 33,07& 28,1 & 35,81  & 33,8     \\
 \rowcolor[gray]{.90}     \textit{Das Zafche -- göcktig} & 37,23& 32,4 & 26,46  & 28,8 &  36,31 & 34,1 \\
                            \textit{Der Nümig -- fongig} & 34,84& 30,1 & 28,79& 26,8  & 36,37& 29,4 \\
\rowcolor[gray]{.90}    \textit{Das Tusern -- zarwig} & 29,03& 26,1  & 33,64 & 29,6  &  37,33  & 32,8 \\
                        \textit{Der Jordruck -- kieflich} & 29,15  & 29,2&  32,48& 31,0 &  38,37  & 33,4 \\
  \rowcolor[gray]{.90}    \textit{Die Itse -- rümmig} & 30,33  & 27,7 & 30,94 & 27,4  &  38,73 & 32,9  \\
                        \textit{Der Wonger -- zänsig} & 29,31 & 28,6  &  31,93 & 27,8&  38,76 & 32,5  \\
  \rowcolor[gray]{.90}    \textit{Das Zülugs -- wosplich} & 30,79  & 30,0 & 30,43 & 28,5   & 38,78 & 28,7\\
                        \textit{Der Zegeff -- ralklich} & 28,14 & 27,6 &  32,16 & 29,7 &  39,70  & 32,3\\
 \rowcolor[gray]{.90}  \textit{Das Hensuln -- ontig} & 31,88   & 30,5&  28,13 & 26,0 & 39,99 & 32,8 \\
                        \textit{Der Cobe -- speiglich} & 29,30& 29,1 &  30,45   & 29,8 & 40,25& 35,0  \\
 \rowcolor[gray]{.90}  \textit{Der Gedalst -- wübig} & 27,79  & 29,9 & 31,82 & 29,0  & 40,39& 31,6 \\
                          \textit{Das Zeseet -- bühllich} & 30,94 & 30,1 & 28,19  & 26,4 &  40,87  & 31,6 \\
 \rowcolor[gray]{.90}    \textit{Der Zotzer -- galklich} & 28,16 & 28,4  &  30,30 & 29,8&  41,54  & 32,0 \\
                        \textit{Der Pulpart -- röngig} & 28,01  & 25,7&  27,18 & 26,0 &  44,81  & 29,7\\
\bottomrule
\end{tabularx}
\caption[Mittel der gesetzten Taler der Pseudowörter bei {Aufgabe III} von Experiment 1]{Mittel der gesetzten Taler und Standardabweichungen der Kombination aus Pseudonomen und Pseudoadjektiv bei Aufgabe III -- Einstellungspolarität (Talersumme je Wortpaar = 100; angeordnet nach ansteigendem negativen Skalenwert).}
\label{Pseudos3}
\end{table}

\chapter{{Experimentelle Studie 1 -- Regressionsmodelle}} \label{Regr1}



\section{Aufgabe I -- Skalarität}

\subsection{{Tabelle \ref{LinReg0001}}}
\ea
Formel des vollen Modells zu Aufgabe I -- Skalarität (mit `Expressiv' als Referenzlevel):\smallskip

$ lmer(\text{Skalenwert}\sim(\text{Expr vs. Desk}+\text{Expr vs. Unm)} * \text{Position} + (1 + (\text{Expr vs. Desk} + \text{Expr vs. Unm}) * \text{Position}|\text{Versuchspers.}) + (1|\text{Pseudow.})) $
\z



\ea
Formel des finalen Modells zu Aufgabe I -- Skalarität (mit `Expressiv' als Referenzlevel):\smallskip

$ lmer(\text{Skalenwert}\sim\text{Expr vs. Desk}+\text{Expr vs. Unm} + \text{Position} + (1 + (\text{Expr vs. Desk} + \text{Expr vs. Unm}) * \text{Position}|\text{Versuchspers.}) + (1|\text{Pseudow.})) $
\z


\subsection{{Tabelle \ref{LinReg0010}}}
\ea
Formel des vollen Modells zu Aufgabe I -- Skalarität (mit `Unmarkiert' als Referenzlevel):\smallskip

$ lmer(\text{Skalenwert} \sim (\text{Unm vs. Desk} + \text{Unm vs. Expr)} * \text{Position} + (1 + (\text{Unm vs. Desk} + \text{Unm vs. Expr)} * \text{Position|Versuchspers.}) + (1|\text{Pseudow.})) $
\z

\ea
Formel des finalen Modells zu Aufgabe I -- Skalarität (mit `Unmarkiert' als Referenzlevel):\smallskip

$ lmer(\text{Skalenwert} \sim \text{Unm vs. Desk} + \text{Unm vs. Expr} + \text{Position} + (1 + (\text{Unm vs. Desk} + \text{Unm vs. Expr}) * \text{Position|Versuchspers.}) + (1|\text{Pseudow.})) $
\z

\subsection{{Tabelle \ref{LinRea0100}}}
\ea
Formel des vollen Modells zu Aufgabe I -- Skalarität

(Einfluss von Frequenz mit `Hoch' als Referenzlevel, Persistenz mit `Alt' als Referenzlevel):\smallskip

$ lmer(\text{Skalenwert} \sim (\text{Frequenz + Persistenz}) * \text{Position} + (1 +(\text{Frequenz + Persistenz}) * \text{Position|Versuchspers.}) + (1|\text{Pseudow.})) $
\z

\ea
Formel des finalen Modells zu Aufgabe I -- Skalarität

(Einfluss von Frequenz mit `Hoch' als Referenzlevel, Persistenz mit `Alt' als Referenzlevel):\smallskip

$ lmer(\text{Skalenwert} \sim (\text{Frequenz + Persistenz}) * \text{Position} - \text{Frequenz : Position - Persistenz : Position} + (1 + (\text{Frequenz + Persistenz}) * \text{Position|Versuchspers.}) + (1|\text{Pseudow.})) $
\z



\section{Aufgabe II -- Sprechereinstellung}

\subsection{{Tabelle \ref{LinReg0002}}}
\ea
Formel des vollen Modells zu Aufgabe II -- Sprechereinstellung (mit `Expressiv' als Referenzlevel):\smallskip

$ lmer(\text{Skalenwert} \sim (\text{Expr vs. Desk} + \text{Expr vs. Unm}) * \text{Position} + (1 + (\text{Expr vs. Desk} + \text{Expr vs. Unm}) * \text{Position|Versuchspers.}) + (1|\text{Pseudow.}))$
\z

\ea
Formel des finalen Modells zu Aufgabe II -- Sprechereinstellung (mit `Expressiv' als Referenzlevel):\smallskip

$ lmer(\text{Skalenwert} \sim \text{Expr vs. Desk} + \text{Expr vs. Unm} + (1 + (\text{Expr vs. Desk} + \text{Expr vs. Unm}) * \text{Position|Versuchspers.}) + (1|\text{Pseudow.})) $
\z


\subsection{{Tabelle \ref{LinReg0020}}}

\ea
Formel des vollen Modells zu Aufgabe II -- Sprechereinstellung (mit `Unmarkiert' als Referenzlevel):\smallskip

$ lmer(\text{Skalenwert} \sim (\text{Unm vs. Desk} + \text{Unm vs. Expr}) * \text{Position} + (1 + (\text{Unm vs. Desk} + \text{Unm vs. Expr}) * \text{Position|Versuchspers.}) + (1|\text{Pseudow.})) $
\z

\ea
Formel des finalen Modells zu Aufgabe II -- Sprechereinstellung (mit `Unmarkiert' als Referenzlevel):\smallskip

$ lmer(\text{Skalenwert} \sim \text{Unm vs. Desk} + \text{Unm vs. Expr} + (1 + (\text{Unm vs. Desk} + \text{Unm vs. Expr}) * \text{Position|Versuchspers.}) + (1|\text{Pseudow.})) $
\z

\subsection{{Tabelle \ref{LinRegzaz0}}}
\ea
Formel des vollen Modells zu Aufgabe II -- Sprechereinstellung

(Einfluss von Frequenz mit `Hoch' als Referenzlevel, Persistenz mit `Alt' als Referenzlevel):
\smallskip

$ lmer(\text{Skalenwert} \sim (\text{Frequenz + Persistenz}) * \text{Position}) + (1 +(\text{Frequenz + Persistenz}) * \text{Position|Versuchspers.}) + (1|\text{Pseudow.})) $
\z

\ea Formel des finalen Modells zu Aufgabe II - Sprechereinstellung 

(Einfluss von Frequenz mit Hoch' als Referenzlevel, Persistenz mit Alt' als Referenzlevel):
\smallskip


$ lmer(\text{Skalenwert} \sim (\text{Frequenz + Persistenz}) * \text{Position}) - \text{Position} - \linebreak \text{Persistenz} - \text{Frequenz} : \text{Persistenz} - \text{Persistenz} : \text{Position} - \text{Frequenz} : \text{Position} + (1 + (\text{Frequenz} + \text{Persistenz}) * \text{Position}|\text{Versuchspers.}) + (1|\text{Pseudow.}) $
\z


\subsection{{Tabelle \ref{LM01}}}
\ea Formel des linearen Modells zum Einfluss von Sprechereinstellung auf Skalarität:
\smallskip


$ lm(\text{Skalarität} \sim \text{Sprechereinstellung}) $
\z


\section{Aufgabe III -- Einstellungspolarität}



\subsection{{Tabelle \ref{LinReg0004}}}
\ea

Formel des vollen Modells zu Aufgabe III -- Einstellungspolarität (mit `Neutral' als Referenzlevel):
\smallskip


$ lmer(\text{Neutral} \sim (\text{Sprechereinstellung + Skalarität}) * \text{Position} + (1 + (\text{Sprechereinstellung + Skalarität}) * \text{Position|Versuchspers.}) +  (1|\text{Pseudow.})) $
\z

\ea
Formel des finalen Modells zu Aufgabe III -- Einstellungspolarität (mit `Neutral' als Referenzlevel):
\smallskip


$lmer(\text{Neutral} \sim (\text{Sprechereinstellung + Skalarität}) * \text{Position} - \text{Sprechereinstellung : Skalarität} - \text{Sprechereinstellung : Position} - \text{Skalarität : Position} + (1 + (\text{Sprechereinstellung + Skalarität}) * \text{Position|Versuchspers.}) + (1|\text{Pseudow.})) $
\z

\chapter{{Experimentelle Studie 2a -- Ratings der Adjektive}}   \label{Aja1}
\begin{table}

\begin{tabularx}{\textwidth}{l rY rY rY}
\lsptoprule
  \multicolumn{7}{c}{\textbf{Experiment 2a} \textbf{Skalarität}}  \\
  \midrule
 & \multicolumn{4}{c}{{Intensiviererklasse}}   \\
 \cmidrule{2-5}
{Adjektiv} &  deskriptiv &   \textit{SD} & expressiv &\textit{SD} &  {Mittelwert} &  \textit{SD}\\
\midrule
 \rowcolor[gray]{.90} \textit{breit} & 55,87& 23,5 & 60,25& 24,3 & 58,06 & 24,0 \\
                        \textit{lang} & 54,90  & 23,8& 62,03  & 22,5& 58,47 & 23,4 \\
  \rowcolor[gray]{.90} \textit{dick} & 57,27 & 23,4& 60,32  & 23,4 & 58,80 & 23,4 \\
                      \textit{groß} & 57,84 & 22,3 & 61,32 & 23,9  & 59,58 & 23,2 \\
  \rowcolor[gray]{.90}  \textit{hoch} & 58,33  & 23,8& 63,37 & 22,8 & 60,85 & 23,5  \\
                        \textit{alt} & 60,28 & 23,1 & 64,56 & 24,8  & 62,42& 24,0 \\
  \rowcolor[gray]{.90} \textit{heiß} & 62,79 & 23,6 & 66,82 & 25,1& 64,80 & 24,4 \\
\lspbottomrule

\end{tabularx}
\caption[Gemittelte Skalenwerte der Adjektive beider Intensiviererklassen von Experiment 2a]{Gemittelte Skalenwerte und Standardabweichungen der Adjektive beider Intensiviererklassen von Experiment 2a -- Skalarität (angeordnet nach ansteigendem Mittelwert).}
\label{Adj0}
\end{table}

\chapter{{Experimentelle Studie 2a -- Regressionsmodelle}}   \label{RegX}

\section{Experiment 2a -- Skalarität}

\subsection{{Tabelle \ref{LinReg01}}}

\ea
Formel des vollen Modells zu Experiment 2a -- Skalarität (mit `Deskriptiv' als Referenzlevel):
\smallskip


$ lmer(\text{Skalenwert} \sim (\text{Desk vs. Expr}) * \text{Position} + (1 + (\text{Desk vs. Expr}) * \text{Position|Versuchspers.}) + (1|\text{Adj.})) $
\z

\ea
Formel des finalen Modells zu Experiment 2a -- Skalarität (mit `Deskriptiv' als Referenzlevel):
\smallskip


$ lmer(\text{Skalenwert}\sim(\text{Desk vs. Expr}) * \text{Position} + (1 + (\text{Desk vs. Expr}) *\linebreak \text{Position}|\text{Versuchspers.}))$
\z


\subsection{{Tabelle \ref{Matrix_1_neu}}}

\ea
Formel der Paarvergleiche zu Experiment 2a -- Skalarität:
\smallskip


$ lmer(\text{Skalenwert} \sim \text{Intensivierer * Position} + (1|\text{Versuchspers.}) + (1|\text{Adj.})) $
\z


\chapter{{Experimentelle Studie 2b -- Ratings der Adjektive}}       \label{Aja3}
\begin{table}

\begin{tabularx}{\textwidth}{l rY rY rY}
\lsptoprule
  \multicolumn{7}{c}{\textbf{Experiment 2b Sprechereinstellung}}  \\
  \midrule
 & \multicolumn{4}{c}{{Intensiviererklasse}}   \\
 \cmidrule{2-5}
{Adjektiv} &  deskriptiv &   \textit{SD} & expressiv &\textit{SD} &  {Mittelwert} &  \textit{SD}\\
\midrule
 \rowcolor[gray]{.90} \textit{breit} & 44,75 & 26,3 & 55,13 & 29,7 & 49,94 & 28,5 \\
                      \textit{dick} & 44,60 & 27,4 & 61,25 & 29,1  & 52,93 & 29,4\\
 \rowcolor[gray]{.90} \textit{lang} & 48,54  & 28,4 & 60,68  & 28,6& 54,61  & 29,1\\
                      \textit{groß} & 51,54 & 27,6  & 59,59& 28,3 & 55,57  & 28,2\\
   \rowcolor[gray]{.90} \textit{alt} & 49,44  & 28,8& 62,24 & 28,3  & 55,84 & 29,2\\
                        \textit{hoch} & 50,43 & 28,0 & 61,81  & 27,2 & 56,12 & 28,2\\
  \rowcolor[gray]{.90} \textit{heiß} & 54,73  & 27,3& 62,32 & 28,0  & 58,52 & 27,8\\
\lspbottomrule

\end{tabularx}
\caption[Gemittelte Skalenwerte der Adjektive beider Intensiviererklassen von Experiment 2b]{Gemittelte Skalenwerte und Standardabweichungen der Adjektive beider Intensiviererklassen von Experiment 2b -- Sprechereinstellung (angeordnet nach ansteigendem Mittelwert).}
\label{Adj10}
\end{table}

\chapter{{Experimentelle Studie 2b -- Regressionsmodelle}}   \label{Aja2}


\section{Experiment 2b -- Sprechereinstellung}

\subsection{{Tabelle \ref{LinReg2}}}

\ea
Formel des vollen Modells zu Experiment 2b -- Sprechereinstellung

(mit `Deskriptiv' als Referenzlevel):\\
$ lmer(\text{Skalenwert} \sim (\text{Desk vs. Expr}) * \text{Position} + (1 + (\text{Desk vs. Expr}) * \text{Position|Versuchspers.}) + (1|\text{Adj.}))$
\z

\ea
Formel des finalen Modells zu Experiment 2b -- Sprechereinstellung\\
(mit `Deskriptiv' als Referenzlevel):
\smallskip


$ lmer(\text{Skalenwert} \sim (\text{Desk vs. Expr}) * \text{Position} - \text{Position} - \text{Desk vs. Expr} : \text{Position} + (1 + (\text{Desk vs. Expr}) * \text{Position}|\text{Versuchspers.})) $
\z
\subsection{{Tabelle \ref{LM02}}}

\ea
Formel des linearen Modells zum Einfluss von Sprechereinstellung auf Skalarität:
\smallskip


$lm(\text{Skalarität} \sim \text{Sprechereinstellung})$
\z


\subsection{{Tabelle \ref{Sign2_neu}}}

\ea
Formel der Paarvergleiche zu Experiment 2b -- Sprechereinstellung:
\smallskip

$ lmer(\text{Skalenwert} \sim \text{Intensivierer * Position} + (1|\text{Versuchspers.}) + (1|\text{Adj.}))$
\z
